\documentclass[10pt,a4paper]{amsart}
\usepackage[margin=19mm]{geometry}
\usepackage{amsmath,amssymb,mathtools}
\usepackage[scr=rsfs,cal=euler]{mathalfa}
\usepackage{xfrac}
\usepackage[unicode,hidelinks,bookmarks=false]{hyperref}
\newcommand{\ie}{i.e.\ }
\newcommand{\cf}{cf.\ }
\newcommand{\resp}{respectively\ }
\newcommand{\Subsection}{\subsection}
\providecommand{\nobreakdash}{\nobreak\hbox{-}}
\setlength{\emergencystretch}{2em}
\multlinegap0pt
\usepackage{amssymb}
\usepackage{mathtools}
\newtheorem{theorem}{Theorem}
\newtheorem{lemma}{Lemma}
\title{Sharp mean-width and Jacobian bounds for Euclidean and hyperbolic harmonic maps}
\author{Deguang Zhong, David Kalaj}
\date{Source: arXiv:2609.09316v1 (CC BY 4.0)}
\begin{document}
\maketitle

\noindent\emph{Source and licence.} Sharp mean-width and Jacobian bounds for Euclidean and hyperbolic harmonic maps. Version arXiv:2609.09316v1 (CC BY 4.0), distributed by arXiv under the Creative Commons Attribution 4.0 International licence, http://creativecommons.org/licenses/by/4.0/, as recorded in the arXiv licence metadata for this exact version and shown on the abstract page https://arxiv.org/abs/2609.09316v1. Full licence notice: \texttt{licenses/arxiv-2609.09316-CC-BY-4.0.txt}.\par\smallskip

\noindent\textit{Editorial scope.} Counted unit: the article's theorem thm:zonal-rigidity (source0.tex lines 339--356). The article's complete original proof of it is delivered with it (source0.tex lines 1298--1546, one \texttt{proof} environment of 249 lines, which the article places away from the statement). No mathematics is written, repaired or re-derived: the statement and the proof are the article's own lines, in the article's own notation, and no retained line is substituted or re-flowed.  Retained uncounted as the source-local context the proof needs: lines 316--330; lines 898--914; lines 975--986; lines 1256--1264.  Nothing was omitted from any retained span, so the package carries no omission record.  No retained line contains a citation, so the wrapper carries no bibliography and no bibliography entry.  No retained line contains a figure environment, an image-inclusion command or drawing code, so no image is delivered and the package declares no asset dependency.  Editorial formatting only, in addition: an AMS \texttt{amsart} preamble that restates the article's own declarations and macros used by the retained lines, namely the environments \texttt{theorem}, \texttt{lemma} on separate counters, together with the standard packages those lines need.  The retained environments print in this document as Theorem 1, Lemma 1 in order of appearance, whereas the article prints the same items with the numbering of its own full document.  The complete native source of the article as distributed in the arXiv e-print for this exact version is bundled unchanged as \texttt{sources/209801/source0.tex}.
\begin{theorem}[Sharp monotone zonal contraction]\label{thm:zonal-main}
Let $n\ge2$, let $k:[0,\pi]\to\mathbb R$ be continuous and nonincreasing,
and let
\[
 F:\mathbb S^{n-1}\to\overline{\mathbb B^n}
\]
be measurable.  Then
\begin{equation}\label{eq:zonal-main}
 w\!\left(
 \operatorname{co}T_kF(\mathbb S^{n-1})
 \right)
 \le2\lambda_1(k).
\end{equation}
The constant is sharp.
\end{theorem}

\begin{lemma}[Measurable maximizing selectors]\label{lem:quantile-selector}
Fix $s\in(0,1)$.  There exists a jointly measurable function
\[
 (u,\eta)\longmapsto a_{u,s}(\eta)\in[0,1]
\]
such that
\[
 \int_{\mathbb S^{n-1}} a_{u,s}(\eta)\,d\sigma(\eta)=s
\]
and
\[
 M_u(s)
 =
 \int_{\mathbb S^{n-1}} a_{u,s}(\eta)\langle F(\eta),u\rangle\,d\sigma(\eta)
\]
for every $u$ outside a null set.
\end{lemma}

\begin{theorem}[Sharp trimmed-body mean-width inequality]\label{thm:trimmed}
For every measurable
$F:\mathbb S^{n-1}\to\overline{\mathbb B^n}$ and every $0\le s\le1$,
\begin{equation}\label{eq:trimmedbound}
 \int_{\mathbb S^{n-1}}M_u(s)\,d\sigma(u)\le \Phi_n(s).
\end{equation}
Equivalently,
\[
 w(\mathcal Z_s(F))\le 2\Phi_n(s).
\]
The estimate is sharp.
\end{theorem}

\begin{lemma}[Continuity in the trimming parameter]\label{lem:M-continuous}
For fixed measurable $F$ with $|F|\le1$ and fixed
$u\in\mathbb S^{n-1}$, the function $s\mapsto M_u(s)$ is
$1$-Lipschitz on $[0,1]$.  Consequently
\[
 J(s):=\int_{\mathbb S^{n-1}}M_u(s)\,d\sigma(u)
\]
is continuous.
\end{lemma}

\begin{theorem}[Rigidity for strictly decreasing zonal kernels]
\label{thm:zonal-rigidity}
Under the hypotheses of Theorem~\ref{thm:zonal-main}, assume in addition
that $k$ is strictly decreasing.  Then
\[
 w\!\left(
 \operatorname{co}T_kF(\mathbb S^{n-1})
 \right)
 =
 2\lambda_1(k)
\]
if and only if
\[
 F(\eta)=Q\eta
 \qquad\text{for a.e. }\eta\in\mathbb S^{n-1}
\]
for some $Q\in O(n)$.
\end{theorem}

\begin{proof}[Proof of Theorem~\ref{thm:zonal-rigidity}]
The converse direction has already been proved in the sharpness part of
Theorem~\ref{thm:zonal-main}.  Suppose therefore that
\begin{equation}\label{eq:zonal-rigidity-equality}
 w(K_k)=2\lambda_1(k).
\end{equation}
Because $k$ is strictly decreasing, the Stieltjes measure $\nu_k$ has full
support:
\[
 \nu_k((a,b))>0
 \qquad
 (0\le a<b\le\pi).
\]

Let
\[
 J(s)=\int_{\mathbb S^{n-1}} M_u(s)\,d\sigma(u).
\]
The proof of Theorem~\ref{thm:zonal-main} gives
\[
 \int_{\mathbb S^{n-1}} h_{K_k}(u)\,d\sigma(u)
 \le
 \int_{(0,\pi)}J(A_n(t))\,d\nu_k(t)
 \le
 \int_{(0,\pi)}\Phi_n(A_n(t))\,d\nu_k(t)
 =
 \lambda_1(k).
\]
By \eqref{eq:zonal-rigidity-equality}, both inequalities are equalities.
The function
\[
 t\longmapsto
 \Phi_n(A_n(t))-J(A_n(t))
\]
is continuous and nonnegative by Lemma~\ref{lem:M-continuous}.  Since
$\nu_k$ has full support, its integral can vanish only if
\begin{equation}\label{eq:J-equality-all}
 J(s)=\Phi_n(s),
 \qquad 0\le s\le1.
\end{equation}

Fix $s\in(0,1)$ and choose the maximizing selectors of
Lemma~\ref{lem:quantile-selector}.  Define
\[
 q_s(\eta)
 =
 \int_{\mathbb S^{n-1}}a_{u,s}(\eta)\,d\sigma(u).
\]
Then $\int_{\mathbb S^{n-1}} q_s(\eta)\,d\sigma(\eta)=s$.  Retaining all inequalities in the proof of
Theorem~\ref{thm:trimmed}, and using \eqref{eq:J-equality-all}, gives
\begin{align}
 \Phi_n(s)
 &=
 J(s)\notag\\
 &\le
 \int_{\mathbb S^{n-1}} |F(\eta)|\Phi_n(q_s(\eta))\,d\sigma(\eta)\notag\\
 &\le
 \int_{\mathbb S^{n-1}} \Phi_n(q_s(\eta))\,d\sigma(\eta)\notag\\
 &\le
 \Phi_n\!\left(\int_{\mathbb S^{n-1}} q_s(\eta)\,d\sigma(\eta)\right)
 =
 \Phi_n(s).
\label{eq:zonal-rigidity-chain}
\end{align}
Thus equality holds at every stage.  Since
\[
 \Phi_n'(s)=\cos\alpha_n(s)
\]
is strictly decreasing on $(0,1)$, $\Phi_n$ is strictly concave on $[0,1]$.
Equality in Jensen's inequality yields
\begin{equation}\label{eq:zonal-q}
 q_s(\eta)=s
 \qquad\text{for a.e. }\eta.
\end{equation}
As $\Phi_n(s)>0$, equality in the preceding inequality also gives
\begin{equation}\label{eq:zonal-F-sphere}
 |F(\eta)|=1
 \qquad\text{for a.e. }\eta.
\end{equation}

Put
\[
 c_s=\cos\alpha_n(s).
\]
For fixed $v\in\mathbb S^{n-1}$, the unique bathtub maximizer of mass $s$
for the function $u\mapsto\langle u,v\rangle$ is, up to a null set,
\[
 \mathbf1_{\{\langle u,v\rangle>c_s\}}.
\]
Indeed, if $b_*$ denotes this indicator and $b$ is any competitor of mass
$s$, then
\[
 \int_{\mathbb S^{n-1}}(b_*(u)-b(u))(\langle u,v\rangle-c_s)\,d\sigma(u)\ge0,
\]
and equality forces $b=b_*$ away from the zero-measure level set
$\{\langle u,v\rangle=c_s\}$.

Equality in the first step of \eqref{eq:zonal-rigidity-chain}, together
with \eqref{eq:zonal-q} and \eqref{eq:zonal-F-sphere}, therefore gives
\begin{equation}\label{eq:zonal-selector-cap}
 a_{u,s}(\eta)
 =
 \mathbf1_{\{\langle F(\eta),u\rangle>c_s\}}
 \qquad\text{for a.e. }(u,\eta).
\end{equation}
Choose a countable dense set $S_0\subset(0,1)$.  Since every selector has
mass $s$, Fubini and \eqref{eq:zonal-selector-cap} imply that for almost
every $u$ and every $s\in S_0$,
\begin{equation}\label{eq:zonal-tail-dense}
 \sigma\{
 \eta:\langle F(\eta),u\rangle>c_s
 \}
 =
 s.
\end{equation}
The numbers $c_s$, $s\in S_0$, are dense in $(-1,1)$.  For the model
variable $Z(\eta)=\eta_n$ one has
\[
 \sigma\{Z>c_s\}=s,
 \qquad 0<s<1,
\]
by the definition of $c_s=\cos\alpha_n(s)$.  Hence, for almost every $u$,
\eqref{eq:zonal-tail-dense} identifies the tail function
$c\mapsto\sigma\{\langle F(\eta),u\rangle>c\}$ with the tail function of
$Z$ on the dense set $\{c_s:s\in S_0\}$.  Both tail functions are monotone
and right-continuous, while the spherical-coordinate distribution is
continuous.  It follows that the two tail functions agree for every
$c\in(-1,1)$, and therefore
\begin{equation}\label{eq:zonal-projection-law}
 \langle F(\eta),u\rangle
 \quad\text{has the same distribution as}\quad
 \eta_n.
\end{equation}
In particular this distribution has no atoms.

We next return to the first inequality in the proof of
Theorem~\ref{thm:zonal-main}.  Define
\[
 B_k(u)
 =
 k(\pi)\int_{\mathbb S^{n-1}}X_u\,d\sigma
 +
 \int_{(0,\pi)}M_u(A_n(t))\,d\nu_k(t).
\]
Then
\[
 h_{K_k}(u)\le B_k(u)
 \qquad\text{for every }u.
\]
Because \eqref{eq:J-equality-all} holds, integration in $u$ gives
\[
 \int_{\mathbb S^{n-1}}B_k(u)\,d\sigma(u)
 =
 \lambda_1(k)
 =
 \int_{\mathbb S^{n-1}}h_{K_k}(u)\,d\sigma(u).
\]
Hence the nonnegative deficit $B_k-h_{K_k}$ vanishes for almost every
$u$.

Fix such a direction $u$ for which
\eqref{eq:zonal-projection-law} also holds.  Choose
$\xi_u\in\mathbb S^{n-1}$ at which
\[
 \xi\longmapsto
 \int_{\mathbb S^{n-1}} k(d(\xi,\eta))X_u(\eta)\,d\sigma(\eta)
\]
attains its maximum.  Applying the layer-cake identity at this maximizing
direction yields
\[
 B_k(u)-h_{K_k}(u)
 =
 \int_{(0,\pi)}D_u(t)\,d\nu_k(t),
\]
where
\[
 D_u(t)
 =
 M_u(A_n(t))
 -
 \int_{C(\xi_u,t)}X_u(\eta)\,d\sigma(\eta)
 \ge0.
\]
Thus the integral of $D_u$ is zero.  The function $D_u$ is continuous:
the first term is continuous by Lemma~\ref{lem:M-continuous}, while the
second is continuous because spherical caps vary continuously in measure
and $X_u$ is bounded.  Since $\nu_k$ has full support,
\[
 D_u(t)=0,\qquad0<t<\pi.
\]
Thus every cap $C(\xi_u,t)$ is a bathtub optimizer of mass $A_n(t)$ for
$X_u$.  By \eqref{eq:zonal-projection-law},
\[
 \sigma\{X_u>\cos t\}=A_n(t),
\]
and the level set $\{X_u=\cos t\}$ has measure zero.  Hence the unique
optimizer of this mass is $\{X_u>\cos t\}$, and therefore, for every
$t\in(0,\pi)$,
\begin{equation}\label{eq:zonal-superlevel}
 \mathbf1_{\{X_u(\eta)>\cos t\}}
 =
 \mathbf1_{\{\langle\eta,\xi_u\rangle>\cos t\}}
 \qquad\text{for a.e. }\eta.
\end{equation}

For each fixed $t$, the two indicators in
\eqref{eq:zonal-superlevel} agree for almost every $\eta$.  Integrating
their absolute difference over $(0,\pi)\times\mathbb S^{n-1}$ and applying
Tonelli's theorem therefore shows that, for almost every $\eta$, they agree
for almost every $t$.  If two numbers $a,b\in[-1,1]$ satisfy
\[
 \mathbf1_{\{a>\cos t\}}=\mathbf1_{\{b>\cos t\}}
 \quad\text{for a.e. }t\in(0,\pi),
\]
then $a=b$, since otherwise the two indicators differ on a nonempty interval
of $t$-values.  Consequently
\[
 X_u(\eta)=\langle\eta,\xi_u\rangle
 \qquad\text{for a.e. }\eta.
\]
Thus
\[
 \langle F(\cdot),u\rangle
 \in\mathcal H_1(\mathbb S^{n-1})
\]
for almost every $u$.

Finally set
\[
 E=
 \{u\in\mathbb R^n:
 \langle F(\cdot),u\rangle\in\mathcal H_1(\mathbb S^{n-1})\}.
\]
This is a linear subspace of $\mathbb R^n$, and
$E\cap\mathbb S^{n-1}$ has full spherical measure.  Hence $E=\mathbb R^n$.
Therefore every coordinate of $F$ is a first spherical harmonic, so
\[
 F(\eta)=A\eta
 \qquad\text{a.e.}
\]
for some real matrix $A$.  By \eqref{eq:zonal-F-sphere},
$|A\eta|=1$ for almost every $\eta\in\mathbb S^{n-1}$.  Continuity then
gives
\[
 \eta^T(A^TA-I)\eta=0
 \qquad\text{for every }\eta\in\mathbb S^{n-1},
\]
hence $A^TA=I$.  Thus $A=Q\in O(n)$, completing the proof.
\end{proof}

\end{document}
